% Independent contribution. AI-produced by GPT-6 Astra; no priority claimed.
\subsection{The first ten-point square and the harmonic half-character}

The first square of (A.T23) contains information in addition to
its comparison with the second square.  We retain that square,
including its actual peripheral markings.  Throughout this
subsection $F=(1,f)\in\widehat{GT}$ has multiplier one.  Put
\[
 B=\langle s,r\mid s^2=r^3=1\rangle^{\wedge},\quad
 \tau=sr,\quad x=\tau^{-2},\quad y=sxs,\quad z=(xy)^{-1}.
 \tag{A.J1}
\]
The kernel of $B\to\mathfrak S_3$ is the named
$\widehat F(x,y)$.  Use precisely the normalized modular action
and pure harmonic frame supplied in R1:
\[
 \begin{gathered}
 \alpha(\tau)=\tau,\quad \alpha(x)=x,\quad
 \alpha(y)=f(x,y)^{-1}yf(x,y),\\
 \alpha(s)=\beta s\beta^{-1}=sf,
 \qquad (\beta^{-1})^{\rm ab}=(\rho,\rho).
 \end{gathered}
 \tag{A.J2}
\]
These equalities specify the frame $\beta$ used below; we do not
replace it by a primitive chosen from a two-cycle relation.

Define $b=c_2(f)$ integrally as follows.  The index-two subgroup
$\ker(\widehat F(x,y)\to C_2)$, with $x\mapsto1,y\mapsto0$,
has basis $x^2,y,xyx^{-1}$.  The image of $f\in\widehat F'$ in
its abelianization is $(0,b,-b)$.  This definition uses only an
actual open subgroup and its abelianization.

\begin{result}{A.5.1}{Proved}{The harmonic primitive with the specified modular normalization}
For every multiplier-one full profinite GT element,
\[
 c_2(f)=\rho,\qquad
 \boxed{\ \beta^{-1}=f(x,s)x^{2\rho}
                      =f(x,s)(\tau^{-1})^{4\rho}.\ }
 \tag{A.J3}
\]
In the exact notation of Nakamura--Tsunogai, this is their
harmonic primitive relation $(\mathrm{GF}_1)$
\cite[Theorem 2, p.~199, and the comparison on p.~200]{CA_NT2002}.
Its validity here is proved by the ordinary genus-zero tower
square; the cited arithmetic theorem is not used to extend its
domain.
\end{result}
\begin{proof}
The actual first quotient map of the ten-point construction is
\[
 q_b:\widehat F(v,d)\twoheadrightarrow
 R=\langle v,d\mid(vd)^2=1\rangle^{\wedge},\qquad
 K=\widehat F(v,d^2)\triangleleft R.
 \tag{A.J4}
\]
The ordinary tower proof above gives an equality up to one
inner automorphism between its normalized source action and
the target action
\[
 \beta_R(v)=v,\qquad
 \beta_R(d)=f(v,d^2)^{-1}df(v,d^2).
 \tag{A.J5}
\]
In particular its source and target parameters are the same
named $f$, by the cross-ratio and pure-subgroup comparisons in
(A.T13), (A.T18), and (A.T21)--(A.T23).  No higher-genus tower
extension occurs in this step.

We now give the exact change to the harmonic marking.  There
is an isomorphism onto the actual index-three subgroup of $B$,
\[
 j:R\xrightarrow{\sim}R_H=\langle x,s\rangle,
 \qquad j(v)=x,\quad j(d)=x^{-1}s.
 \tag{A.J6}
\]
Both sides have the free-product presentation
$\widehat{\mathbb Z}\amalg C_2$, with respective generators
$(v,vd)$ and $(x,s)$.  In particular
\[
 j(d^2)=x^{-1}y^{-1}=x^{-1}zx.
 \tag{A.J7}
\]
This conjugation is essential to the comparison.  Direct
calculation in $C_2*C_3$ gives
\[
 \theta:=\operatorname{Ad}_{\tau}|_{\widehat F(x,y)}:
 x\mapsto x,\quad y\mapsto x^{-1}y^{-1},
 \qquad
 \theta':=\operatorname{Ad}_x\circ\theta:
 x\mapsto x,\quad y\mapsto (xy)^{-1}.
 \tag{A.J8}
\]
For example $\operatorname{Ad}_s$ exchanges $x,y$ and
$\operatorname{Ad}_r$ cycles $x,y,z$, which proves the first
formula; the second follows by conjugating by $x$.  Since the
genuine normalized action fixes both $\tau$ and $x$, it commutes
literally with both $\theta$ and $\theta'$.  Therefore
$j\beta_Rj^{-1}$ agrees with $\alpha$ on the actual open
subgroup $\widehat F(x,y)=j(K)$, in its basis
$(x,x^{-1}y^{-1})$.  It agrees on $R_H$ as well: the
discrepancy of any extension at an element of $R_H$ centralizes
that open free subgroup, and
$C_{R_H}(\widehat F(x,y))=1$.  The latter follows by intersecting
the procyclic centralizers of the independent peripheral
generators $x,y$ in the free profinite group, or directly from
the free-product action and the nontrivial action of $s$ on
their homology.  Thus the target in this comparison is exactly
the action and the frame of (A.J2).

On the source of (A.J4) choose the free basis
$X=v$, $Y=(vd)^{-1}$.  This is the change $\theta'$ of (A.J8)
in the abstract copy $\widehat F(v,d)$.  The preceding literal
commutation shows that its normalized action is still
$X\mapsto X$, $Y\mapsto f(X,Y)^{-1}Yf(X,Y)$, with no
residual inner factor from the change of basis.  The composite
$j q_b$ is now the actual harmonic opening map
\[
 q_H:\widehat F(X,Y)\twoheadrightarrow R_H,
 \qquad X\mapsto x,\quad Y\mapsto s.
 \tag{A.J9}
\]
Consequently there is $k\in R_H$ such that
$q_H\alpha_{\rm source}=\operatorname{Inn}(k)\alpha q_H$.
Comparison at $X$ gives $k=x^t$ for a unique
$t\in\widehat{\mathbb Z}$, since $C_{R_H}(x)=\langle x\rangle$.
Comparison at $Y$ gives
\[
 f(x,s)^{-1}sf(x,s)=x^t\beta s\beta^{-1}x^{-t}.
\]
It follows that $f(x,s)x^t\beta\in C_{R_H}(s)=\langle s\rangle$.
Every factor on the left is pure: for the first one this also
follows because $R_H/\widehat F(x,y)=C_2$ and $f$ is a
commutator word.  The intersection of this pure subgroup with
$\langle s\rangle$ is trivial.  Hence
\[
                   \beta^{-1}=f(x,s)x^t.                 \tag{A.J10}
\]

The two-cycle gives the abelianization of $f$ in
$\ker(y\bmod2)$, in the basis $(x,y^2,yxy^{-1})$, as
$(-b,0,b)$.  The restriction of (A.J9) sends that basis to
$(x,1,y)$.  Thus $f(x,s)^{\rm ab}=(-b,b)$ in the actual pure
target.  Abelianizing (A.J10) and using (A.J2) yields
$(\rho,\rho)=(-b+t,b)$, so $b=\rho$ and $t=2\rho$.
This proves (A.J3) with its full integral coefficient and
specified modular frame.
\end{proof}

\paragraph{The corresponding braid relation.}
One can read the same first square before changing its source
basis.  Its inner element then centralizes $v$, and comparison
at $d$ gives
\[
 f(v,d)=d^{-t}f(v,d^2)v^{-t}.                             \tag{A.J11}
\]
Here the centralizer of $d$ is $\langle d\rangle$; passing to
$R^{\rm ab}=\widehat{\mathbb Z}[d]\oplus C_2[vd]$ shows
that the two exponents must both equal $-t$ and that $t$ is
even.  In $K^{\rm ab}$ the two sides are respectively
$(-2b,-b)$ and $(-t,-t/2)$: use
$dvd^{-1}=v^{-1}d^{-2}$ and the same index-two Schreier
coefficients.  Hence $t=2b=2\rho$.  The two-cycle rewrites
(A.J11) as
\[
                f(d,v)=v^{2\rho}f(d^2,v)d^{2\rho}.       \tag{A.J12}
\]
In the mixed subgroup $\langle u,w^2\rangle$ of the ordinary
braid group $B_3=\langle u,w\mid uwu=wuw\rangle$, its central
element is $(uw)^3=(uw^2)^2$.  Quotienting by this center gives
(A.J4), with $d=u$, $v=w^2$.  The full completed central kernel
is procyclic: the usual pure subgroup has full topology and is
$F_2\times\mathbb Z$, and is open in this mixed group.
Both evaluations of $f$ have exponent sum zero.  The exponent
sum of the right side of (A.J12) is $6\rho$, while the central
generator has exponent sum six.  Thus the unique central
correction is its $-\rho$ power, giving the exact equality
\[
 f(u,w^2)=w^{4\rho}f(u^2,w^2)u^{2\rho}(uw^2)^{-2\rho}
                         \quad\hbox{in }\widehat B_3.    \tag{A.J13}
\]
This is the first displayed equality of $(\mathrm{IV}')$ in
\cite[Theorem 2.2]{CA_NS2000}, with its coefficient now
identified by (A.J3).  We use (A.J13) only in this explicit
marking and do not infer an extra multiplication relation by
renaming a coefficient.  There is also an actual continuous
surjection
\[
 R\twoheadrightarrow B,\qquad v\mapsto r,\quad
 d\mapsto a=\tau^{-1},
\]
because $ra=s$ and hence $(ra)^2=1$.  Applying this homomorphism
directly to (A.J12), and using $r^3=1$, proves
\[
          f(a,r)=r^{-\rho}f(a^2,r)a^{2\rho}.             \tag{A.J20}
\]
Thus the same first square proves the multiplier-one
central-quotient form of NT's $(\mathrm{IV}'_{\rm bis})$,
through this specified morphism.

\begin{result}{A.5.2}{Proved}{Both normalized parity defects vanish on full GT}
On the multiplier-one domain of the normalized R1 frames,
let $\Psi_3$ be the index-three Schreier character of
DELTA\_ON\_GT A.88 and set
\[
 d_2=\epsilon_2(W)+\Psi_3/2\pmod2,\qquad
 e_3=\kappa_d+\Psi_3/2\pmod2.
 \tag{A.J14}
\]
Then every full GT element in this domain satisfies $d_2=0$.
On the domain where both named frames in (A.J14) are defined,
it also satisfies $e_3=0$.  Thus the individual parity
comparison requested in A.5 holds, in addition to their sum
$\Delta=d_2+e_3$.
\end{result}
\begin{proof}
We include the finite specialization instead of importing a
gamma duplication theorem with additional hypotheses.  Put
$R_2=\mathbb Z_2[\omega]$, $\omega^2+\omega+1=0$, and use
the affine law $(t,u)(t',u')=(t+ut',uu')$.  Evaluate at
$X=(0,\omega)$, $Y=(1,1)$ and write
$f(X,Y)=(t_f,1)$.  The index-three basis
$x^3,y_j=x^jyx^{-j}$ gives
\[
 t_f=m_0+m_1\omega+m_2\omega^2,\quad \sum m_j=0,
 \Psi_{3,2}=m_0=\tfrac13\operatorname{Tr}(t_f).
 \tag{A.J15}
\]
All evaluations are continuous into the profinite affine
group; thus these formulas apply to the actual profinite word.

For completeness the gentle relations imply
$t_f\in\mathbb Z_2$ without using the pentagon.  Their genuine
normalized action has the third-peripheral formula
\[
 \alpha(z)=f(z,x)zf(z,x)^{-1},\qquad z=(xy)^{-1}.           \tag{A.J16}
\]
It also follows immediately from commutation with $\theta'$
and the two-cycle.  At any affine pair with linear components
$U,U^{-1}$, where $U=\omega$ or $\omega^{-1}$, the product
is a translation and its conjugator in (A.J16) is a
translation.  Thus the product is fixed.  If the evaluated
$f$ is the translation $v$, direct multiplication instead
changes that product by $(1-U)v$, so $v=0$.
Return to $X,Y$ and $Z=(XY)^{-1}=(-1,\omega^{-1})$.
The first factor in
$f(Z,X)f(Y,Z)f(X,Y)=1$ is therefore one.  A translation
conjugates $Z$ to $(0,\omega^{-1})$ and leaves $Y$ and all
translation-valued commutator evaluations unchanged.  The
two-cycle consequently makes the second factor the
translation $-\overline{t_f}$.  Hence $t_f=\overline{t_f}$.
The fixed ring is $\mathbb Z_2$, and (A.J15) gives
\[
 \Psi_{3,2}=2t_f/3\in2\mathbb Z_2,\qquad
          \Psi_3/2\bmod2=t_f\bmod2.                    \tag{A.J17}
\]
The odd components have unique halves because two is a unit;
this proves the asserted integral half in the full coefficient
ring as well.

Let $c=(1-\omega)^{-1}\in\mathbb F_4$.  The reduction of
the actual modular affine representation is
\[
 s\mapsto(c,1),\quad r\mapsto(c,\omega),\quad
 \tau\mapsto(0,\omega),\quad
 x\mapsto(0,\omega),\quad y\mapsto(1,\omega).
 \tag{A.J18}
\]
This is precisely the projective level-three map
$B\to\operatorname{PSL}_2(\mathbb F_3)
\cong\mathbb F_4\rtimes C_3$.  One explicit identification
uses the matrices
$S=\left(\begin{smallmatrix}0&-1\\1&0\end{smallmatrix}\right)$,
$R_0=\left(\begin{smallmatrix}0&-1\\1&-1\end{smallmatrix}\right)$:
the bijection from $\mathbb P^1(\mathbb F_3)$ to $\mathbb F_4$
is $0\mapsto\omega$, $1\mapsto1$, $2\mapsto\omega^2$,
$\infty\mapsto0$.  Under it $S$ is $t\mapsto t+c$ and
$R_0$ is $t\mapsto\omega t+c$, with $c=\omega$.
If the simultaneously conjugated matrix convention of the
main text is used, compose this displayed bijection with that
same conjugation; the generator assignments (A.J18) are the
specified ones.  In particular the nontrivial projective
element of the signed centralizer $\langle S\rangle$ is the
translation by $c$.

The harmonic identity (A.J3) now says that the translation
of $g=\beta^{-1}$ is $ct_f$: evaluation at
$((0,\omega),(c,1))$ scales the translation of $f(X,Y)$ by
$c$, and the rightmost power of $x$ has zero translation.
The correcting shear $V=U^{8\rho}$, $U=\tau^{-1}$, also has
zero translation.  The R1 frame $W=gV$ belongs to the signed
centralizer of $S$ at level three, so its projective image is
one or the translation by $c$, and its exponent is exactly
$\epsilon_2(W)$.  Consequently
\[
       \epsilon_2(W)=t_f=\Psi_3/2\pmod2,                \tag{A.J19}
\]
which proves $d_2=0$.  Finally the CM comparison already proved
in A.7 gives $\Delta=0$ on full GT, and the exact identity
$\Delta=d_2+e_3$ gives $e_3=0$ on its stated domain.  This last
step uses that proved comparison, without assuming an
equianharmonic opening identity or identifying an unmarked
Kummer character.
\end{proof}
